% ==============================================================================
% Chapter 8: Multi-Agent Development Methodology & Red Team Audit
% ==============================================================================
\chapter{מתודולוגיית הפיתוח: ועדת המומחים של הסוכנים וביקורת איכות}
\label{ch:methodology}

\section{יישום מתודולוגיית שיעור 4 להפקת דוח מכין רב-תחומי}
מעבר להיותו של פרויקט \en{CivicShield AI} תכנון למערכת רב-סוכנית, תהליך הכתיבה, המידול, הסימולציה וההידור של דוח מכין זה עצמו בוצע במודל \textbf{ועדת מומחים מבוזרת של סוכני קוד ומחקר (\en{Multi-Agent Engineering Committee})}, ברוח עקרונות שיעור 4 בקורס ("תזמור רב-סוכנים עם סוכני קידוד") \cite{segal2026l04}.

במקום לבקש מסוכן יחיד לכתוב ברצף עשרות עמודים בעולמות תוכן זרים זה לזה (משפט חוקתי, אפידמיולוגיה מתמטית, סוציולוגיה של התקשורת ועיצוב \en{LuaLaTeX}), פוצל הפרויקט ללוח משימות מובנה (\en{team.json} ו-5 קובצי הנחיית תפקיד \en{agents/*/AGY.md}), שבו כל סוכן מחזיק בזהות אפיסטמית, כלים ייעודיים וסמכות בקרה הדדית:
\begin{itemize}
    \item \textbf{\en{Master Orchestrator} (\en{agents/master/AGY.md}):} אחראי על ניהול מסמך הדרישות (\en{PRD.md}), תוכנית הביצוע (\en{PLAN.md}), עקביות המספור בין פרקים ואיחוד השכבות לכדי קובץ \en{main.tex} אחיד.
    \item \textbf{\en{Political \& Communication Theory Agent} (\en{agents/theory/AGY.md}):} מתמחה בספרות התקשורת הפוליטית ומדעי המדינה, אחראי על פרקים 2 ו-6 ועל המיפוי ל-26 המקורות הביבליוגרפיים.
    \item \textbf{\en{System \& Security Architect Agent} (\en{agents/architect/AGY.md}):} מתמחה בארכיטקטורת מערכות, סכימת העברת ההודעות (\en{Hub JSON}), פרוטוקולי קריפטוגרפיה ותכנון תרשימי ה-\en{TikZ}.
    \item \textbf{\en{Quantitative Simulator Agent} (\en{agents/simulator/AGY.md}):} פועל בסביבת \en{Python 3 / uv}, כותב ומריץ סימולציות דטרמיניסטיות ב-\en{numpy} ו-\en{matplotlib} להפקת ששת הגרפים המדעיים ומזין את התוצאות המאומתות לפרקים 3, 5 ו-7.
    \item \textbf{\en{Red Team \& RTL Audit Agent} (\en{agents/redteam/AGY.md}):} סוכן לעומתי שתפקידו לתקוף הנחות מוצא אידאולוגיות או הנדסיות, לוודא היעדר צנזורה אוטומטית ולאמת את שלמות הקימפול הדו-כיווני ב-\en{LuaLaTeX}.
\end{itemize}

איור~\ref{fig:committee_workflow} מציג את תרשים הזרימה של ועדת המומחים, כולל שערי החזרה לתיקון (\en{Rejection Feedback Loop}) כאשר הצוות האדום מזהה פגם מתמטי, משפטי או טיפוגרפי.

\begin{figure}[htbp]
\begin{tikzhebrew}
\begin{tikzpicture}[
    node distance=0.8cm and 0.7cm,
    box/.style={rectangle, rounded corners=4pt, draw=CivicNavy, fill=CivicGrayBg, line width=1pt, minimum width=3.4cm, minimum height=1.25cm, align=center, font=\scriptsize},
    orch/.style={rectangle, rounded corners=4pt, draw=CivicNavy, fill=CivicNavy, text=white, line width=1.2pt, minimum width=4.6cm, minimum height=1.25cm, align=center, font=\small},
    red/.style={rectangle, rounded corners=4pt, draw=CivicRed, fill=CivicRedBg, line width=1.2pt, minimum width=4.8cm, minimum height=1.25cm, align=center, font=\footnotesize},
    ok/.style={rectangle, rounded corners=4pt, draw=CivicEmerald, fill=CivicTealBg, line width=1.3pt, minimum width=4.4cm, minimum height=1.2cm, align=center, font=\footnotesize},
    arr/.style={-{Stealth[length=2.3mm]}, line width=1.05pt, color=CivicNavy}
]
    \node[orch] (M) at (0, 4.5) {\textbf{\he{סוכן \en{Master}: אפיון \en{PRD} והזנת לוח משימות}}\\[2pt]{\tiny\en{team.json + TODO Dispatch}}};

    \node[box] (T) at (-4.3, 2.3) {\textbf{\he{סוכן תיאוריה ומשפט}}\\[2pt]{\tiny\he{פרקים 2 ו-6 + ביבליוגרפיה}}};
    \node[box] (A) at (0.0, 2.3) {\textbf{\he{סוכן ארכיטקטורה}}\\[2pt]{\tiny\he{פרקים 1, 3, 4 + תרשימי \en{TikZ}}}};
    \node[box] (S) at (4.3, 2.3) {\textbf{\he{סוכן סימולציות \en{Python}}}\\[2pt]{\tiny\he{פרקים 5, 7 + 6 גרפי \en{PDF/PNG}}}};

    \node[red] (R) at (0, 0.1) {\textbf{\he{שער ביקורת: צוות אדום ובדיקת \en{LuaLaTeX / BiDi}}}\\[2pt]{\tiny\he{Red Team Adversarial Audit + \en{verify\_report.py}}}};

    \node[ok] (F) at (0, -2.0) {\textbf{\he{הפקה סופית מאושרת (\en{CivicShield\_AI\_Preliminary\_Report.pdf})}}\\[2pt]{\tiny\he{אפס שגיאות קימפול, דיוק מחקרי ודחיפה ל-\en{GitHub}}}};

    \draw[arr] (M.south) -- ++(0,-0.45) -| (T.north);
    \draw[arr] (M.south) -- (A.north);
    \draw[arr] (M.south) -- ++(0,-0.45) -| (S.north);

    \draw[arr] (T.south) -- ++(0,-0.45) -| (R.north);
    \draw[arr] (A.south) -- (R.north);
    \draw[arr] (S.south) -- ++(0,-0.45) -| (R.north);

    \draw[arr, color=CivicEmerald] (R.south) -- node[right, font=\scriptsize]{\he{עבר אימות מלא}} (F.north);

    % Rejection feedback loop
    \draw[arr, dashed, color=CivicRed] (R.west) -- ++(-4.5,0) |- node[left, pos=0.25, font=\scriptsize, align=center]{\he{דחייה לתיקון}\\\en{(Fix Loop)}} (M.west);
\end{tikzpicture}
\end{tikzhebrew}
\caption{זרימת הפקת הדוח המכין באמצעות ועדת חמשת הסוכנים ומעגל התיקונים של הצוות האדום}
\label{fig:committee_workflow}
\end{figure}

\section{תצורת צוות הסוכנים ומנגנון בקרת איכות (\en{team.json} ו-\en{AGY.md})}
בהתאמה למבנה הפרויקט בסביבת \en{Antigravity}, מוגדרת המערכת ברמת תיקיית השורש על ידי הקובץ \en{team.json} והקובץ \en{AGY.md}, ומתפצלת לחמש תת-תיקיות סוכנים. מבנה זה מאפשר הרצה מחזורית שבה ניתן להחליף או לשדרג סוכן בודד (למשל הרצת תרחישי קיטוב חדשים אצל סוכן ה-\en{Simulator}) מבלי לשבור את ניסוח הפרקים המשפטיים או התיאורטיים. טבלה~\ref{tab:committee_roster} מרכזת את חלוקת הבעלות הקובצית בוועדה.

{\renewcommand{\arraystretch}{1.35}
\begin{table}[htbp]
\centering
\caption{סגל ועדת הסוכנים המבצעית של פרויקט CivicShield AI וחלוקת הבעלות בקבצים}
\label{tab:committee_roster}
\begin{tabularx}{\textwidth}{|p{2.1cm}|p{4.0cm}|p{3.5cm}|X|}
\hline
\textbf{\textcolor{CivicNavy}{מזהה סוכן}} & \textbf{\textcolor{CivicNavy}{קובץ הנחיות תפקיד}} & \textbf{\textcolor{CivicNavy}{תוצרים בבעלותו}} & \textbf{\textcolor{CivicNavy}{סמכות בקרה מיוחדת}} \\ \hline
\en{master} & {\small\en{agents/master/AGY.md}} & {\small\en{PRD.md}, \en{PLAN.md}, \en{main.tex}} & תזמור לוח משימות והכרה בעמידה ביעדי אורך ומסירה. \\ \hline
\en{theory} & {\small\en{agents/theory/AGY.md}} & פרקים 2, 6 + ביבליוגרפיה & וידוא עקביות תיאורטית ודיוק משפט חוקתי ופרטיות. \\ \hline
\en{architect} & {\small\en{agents/architect/AGY.md}} & פרקים 1, 3, 4, {\small\en{civicshield.cls}} & פיקוח על הפרדת שכבות, תרשימי \en{TikZ} וסכמת \en{Hub}. \\ \hline
\en{simulator} & {\small\en{agents/simulator/AGY.md}} & {\small\en{generate\_figures.py}}, פרקים 5, 7 & וידוא כי כל מספר בדוח נגזר מהרצה מאומתת של הקוד. \\ \hline
\en{redteam} & {\small\en{agents/redteam/AGY.md}} & {\small\en{verify\_report.py}}, פרק 8 & זכות חסימת מסירה במקרה של פגיעה בזכויות או כשל \en{LaTeX}. \\ \hline
\end{tabularx}
\end{table}}

\section{תיעוד ממצאי הצוות האדום והתיקונים שבוצעו בדוח}
יעילותה של גישת \en{Multi-Agent} נמדדת ביכולתה לזהות שגיאות עוד בשלב התכנון המוקדם. במהלך הכנת מהדורה \en{1.0} של הדוח המכין, סיפק סוכן ה-\en{Red Team} ארבע ביקורות משמעותיות שתוקנו באופן מלא טרם האישור הסופי (טבלה~\ref{tab:redteam_log}).

{\renewcommand{\arraystretch}{1.35}
\begin{table}[htbp]
\centering
\caption{יומן ממצאי ביקורת הצוות האדום (Red Team Audit) והתיקונים ההנדסיים שבוצעו}
\label{tab:redteam_log}
\begin{tabularx}{\textwidth}{|p{1.9cm}|p{3.4cm}|p{3.5cm}|X|}
\hline
\textbf{\textcolor{CivicNavy}{ממצא}} & \textbf{\textcolor{CivicNavy}{תיאור הכשל בטיוטה}} & \textbf{\textcolor{CivicNavy}{סיכון מזוהה}} & \textbf{\textcolor{CivicNavy}{תיקון שהוטמע בגרסה הסופית}} \\ \hline
\en{RT-01} & גישור קיטוב שדחס את כל הדעות לאחדות (\en{Uniform Consensus}) & חשש מהנדסת דעות אנטי-פלורליסטית & הוספת עוגן אותנטיות $0.28\,x_i(0)$ למשוואה~\ref{eq:hk_bridge}, המשמר גיוון דמוקרטי מלא. \\ \hline
\en{RT-02} & היתר עקיף לסוכן להמליץ על הסרה בלעדית ללא ערכאה & מדרון חלקלק לצנזורה בידי הרשות המבצעת & הוספת מפתח שלישי (אופוזיציה/ציבור) ב-\en{3-of-3 Ed25519 Quorum} ועץ חוקה קשיח. \\ \hline
\en{RT-03} & גלישת טקסט וסימני כיוון \en{RTL} בקופסאות ובאיורי \en{TikZ} & פגיעה בקריאות שער המכרז ובמשוואות & פיתוח המעטפות \en{\textbackslash NewEnviron} ו-\en{tikzhebrew} במחלקה \en{civicshield.cls} לאיטום כיווניות. \\ \hline
\en{RT-04} & התאמה חלקית בגרפים (\en{fig1, fig5}) ופער בהצלבות תגים & חוסר מיקוד ויזואלי או תווי הגהה & כיול גרפי שיא ו-\en{Erlang-C} בקובץ הפייתון ובניית סקריפט \en{verify\_report.py} לווידוא אוטומטי. \\ \hline
\end{tabularx}
\end{table}}

\section{סיכום הדוח המכין ומעבר לשלב היישום (Handover)}
דוח מכין זה הציב מענה שלם, רב-תחומי ומגובה כמותית לאתגר המרכזי של הדמוקרטיה בעידן הבינה המלאכותית: כיצד להגן על הספרה הציבורית, סדר היום האזרחי וטוהר הבחירות מול מתקפות \en{FIMI} ו-\en{Deepfakes} — מבלי להקריב את חופש הביטוי, הפלורליזם הדעות והפרטיות הפרטית.

באמצעות חיבור של שבע תיאוריות מבוססות בתקשורת פוליטית ומדעי המדינה, ארכיטקטורה מבוזרת של 5 שכבות ותזמור רב-סוכני, ארבעה מודלים מתמטיים מוכחים בסימולציות \en{Python} והגנות חוקתיות קריפטוגרפיות, תשתית \en{CivicShield AI} בשלה למעבר מיידי משלב הדוח המכין אל שלב א' של ההקמה ההנדסית (\en{Months 1--6}). המערכת מדגימה הלכה למעשה כיצד בינה מלאכותית מתוזמרת היטב אינה מאיימת על מוסדות הדמוקרטיה — אלא משמשת חומת מגן שנועדה לשמר אותם שקופים, פתוחים וחסונים לדורות הבאים.
